
\section{SAT Framework — Formal Development Summary: Phase~I and Phase~II}

This document formalizes the foundational stages of the SAT emergent-topological model, comprising \textbf{Phase~I} (Topological Structure and Configuration Space) and \textbf{Phase~II} (Dynamics and Interaction Fields). These stages define the underlying geometry, establish quantization and linking rules, and derive the field equations that govern emergent matter, gauge forces, and time structure.

\tableofcontents

%-------------------------------------------------
\section{Phase I — Structural Geometry and State Space}

\subsection{I.1~ Filament Configuration Space}

\begin{itemize}
  \item \textbf{Topological Base}: Let $M$ be a smooth, connected 4-manifold with no background metric.
  \item \textbf{Filament Set}: Define smooth embeddings $\gamma: \mathbb{R} \rightarrow M$ such that:
  \[
      \mathcal F = \left\{ \gamma \in C^\infty(\mathbb R, M) \;\middle|\; \lim_{\lambda\to\pm\infty} \|\dot{\gamma}^\mu\| < \infty \right\} \Big/ \sim,
  \]
  where $\sim$ denotes ambient isotopy in $M$.
  \item \textbf{Measure}: Equip $\mathcal{F}$ with curvature-weighted Boltzmann measure:
  \[
  \mu[\gamma] \propto \exp\left[-\tfrac{T}{2} \int d\lambda\, \|\ddot\gamma\|^2\right].
  \]
\end{itemize}

\subsection{I.2~ Emergent Geometry Theorem}

Given $F(x)\subset\mathcal F$ (filaments intersecting $x$), define:
\[
  \tilde g_{\mu\nu}(x) = \langle v_\mu v_\nu \rangle_{F(x)},\quad v^\mu = \dot\gamma^\mu.
\]

\paragraph{Theorem (Metric Emergence):}
$\tilde g_{\mu\nu}(x)$ is Lorentzian iff:
\begin{enumerate}
  \item \textbf{Spanning}: filaments span $T^*_xM$
  \item \textbf{Non-degeneracy}: $\det\tilde g_{\mu\nu} \neq 0$
  \item \textbf{Flow Dominance}: $\exists u^\mu$ with $\tilde g_{\mu\nu} u^\mu u^\nu < 0$
\end{enumerate}

Failure yields degenerate or Euclidean behavior—no resolved time.

\subsection{I.3~ Phase Quantization}

Binding occurs when:
\[
   \phi_i - \phi_j = \frac{2\pi k}{n},\quad k \in \mathbb{Z},\; n \in \mathbb{Z}_{>0}
\]
arising from:
\begin{itemize}
  \item Compact holonomy over $T^3$: $\oint d\phi = 2\pi k$
  \item Binding potential: $V_{\text{bind}} \propto -\sum_n \cos[n(\phi_i - \phi_j)]$
\end{itemize}

\subsection{I.4~ Quantum State Space of SAT}

The Hilbert space of SAT is:
\[
    \mathcal H_{\text{SAT}} = L^2(\mathcal C, \mu),\quad \mathcal C = \mathcal F/\sim
\]
Basis states $|\Gamma\rangle$ are labeled by:
- topological charge $Q$,
- winding $w$,
- linking number,
- and phase vectors.

Operators include:
\[
\hat Q|\Gamma\rangle = Q|\Gamma\rangle,\quad \hat m = \frac{m_0}{\hat Q},\quad [\hat\phi,\hat\pi_\phi] = i.
\]

%-------------------------------------------------
\section{Phase II — Field Dynamics and Interaction Rules}

\subsection{Unified Action}

\begin{align}
S_{\text{SAT}} = \int d^4x \sqrt{-g} \Big[ \frac{1}{2\kappa} R
+ \sum_G \frac{1}{4g_G^2} \text{Tr}(F^{(G)}_{\mu\nu} F^{\mu\nu}_{(G)}) 
+ \bar{\Psi}(i\gamma^\mu D_\mu - \tfrac{m_0}{Q})\Psi 
+ \mathcal{L}_{\text{self}} \Big]
\end{align}

\subsection{II.1~ Derived Field Equations}
\begin{align}
& \text{(Gravity)}\quad R_{\mu\nu} - \tfrac{1}{2} R g_{\mu\nu} = \kappa T_{\mu\nu} \\
& \text{(Gauge)}\quad D^\mu F^a_{\mu\nu} = g_G^2 \bar\Psi \gamma_\nu T^a \Psi \\
& \text{(Matter)}\quad (i\gamma^\mu D_\mu - \tfrac{m_0}{Q}) \Psi = 0
\end{align}

\subsection{II.2~ Projective Time and Foliation}

Time emerges from dominant filament alignment:
\begin{itemize}
  \item $u^\mu = v^\mu / \sqrt{-\tilde g_{\alpha\beta}v^\alpha v^\beta}$
  \item $\Sigma_t$: hypersurfaces orthogonal to $u^\mu$
  \item Proper time: $d\tau = \sqrt{-g_{\mu\nu} \dot\gamma^\mu \dot\gamma^\nu} d\lambda$
\end{itemize}

\subsection{II.3~ Topological Origin of Gauge Sectors}

\begin{center}
\begin{tabular}{|c|c|}
\hline
Gauge Group & SAT Topology \\\hline
$U(1)$ & Filament phase twist \\
$SU(2)$ & Hopf-linked bundles \\
$SU(3)$ & Borromean triplets \\\hline
\end{tabular}
\end{center}

Color confinement follows naturally:
- $Q=1$ filaments are unstable alone.
- Only $Q=2,3$ linked states form persistent structures.

%-------------------------------------------------
\section{Conclusion and Forward Trajectory}

Phases I and II establish SAT’s internal structure, projection geometry, and interaction rules. Metric, time, and gauge sectors all emerge from bundle topology and filament dynamics. With this foundation, Phase III maps these structures onto Standard Model particles, and Phase IV expands into cosmology and projective spacetime.

\bigskip
\noindent\textbf{Declared Assumptions:}
A1--A6 (see introduction), covering configuration space, metric conditions, dominant filament flow, and phase locking.
